% day2_seg2_ranproc2/rand_ramp.tex
\begin{frame}
  \onslide<1->{
	\frametitle{Some Important Random Processes} 
	\framesubtitle{Random Constant and Random Ramp}
}
\begin{itemize}
  \onslide<2->{
 \item A random constant process $x$ such that the sample trajectories are constant with respect to time\\
\begin{footnotesize} But the constant value is different for each sample trajectory in the ensemble\\
$x(\omega_1,t)\neq x(\omega_2,t)$, for $\omega_1\neq \omega_2$
\end{footnotesize}
}
  \onslide<3->{
 \item A random ramp requires two state variables: a random constant (y-intercept or bias) and the integral of a
 random constant (yields random slope)
}
  \onslide<4->{
 \item In principle,
 however, this constant value could vary with $\omega$, i.e., across the
 ensemble, thus making the process non-Ergodic.
}
  \onslide<5->{
 \item Random constants are very useful in many applications to model biases that vary from
 turn-on to turn-on
}
\end{itemize}
\end{frame}
\begin{frame}
  \onslide<1->{
	\frametitle{Some Important Random Processes} 
	\framesubtitle{Random Constant and Random Ramp}
}
\begin{itemize}
  \onslide<2->{
 \item See Matlab Simulink demos:  rand\_const.mdl and
 rand\_ramp.mdl
}
  \onslide<3->{
 \item In modeling gyro biases, a combination of a random constant
 (bias stability) and random walk (bias drift) is common.
}
  \onslide<4->{
 \item Gyro scale factors are modeled similarly.
}
  \onslide<5->{
 \item Deterministic biases and scale factors are also part of the model, but these generally are calibrated out.
}
\end{itemize}
\end{frame}
